(1) 忽略边缘效应，高斯定律给 \(E=Q/(\epsilon_0\pi a^2)\)。位移电流密度为
\[
J_d=\epsilon_0\frac{\partial E}{\partial t}
=\frac I{\pi a^2}.
\]
在整个板面积积分得到位移电流 \(I_d=I\)，数值恰为导线中的传导电流。

(2) 由安培–Maxwell 方程
\[
\oint\boldsymbol B\cdot d\boldsymbol\ell
=\mu_0\left(I_{\rm conduction}+\epsilon_0\frac{d\Phi_E}{dt}\right)
\]
对板间半径 r<a 的同轴圆回路计算，穿过圆盘的位移电流为 \(I r^2/a^2\)，所以
\[
2\pi rB_\phi=\mu_0I r^2/a^2,\quad
B_\phi(r)=\frac{\mu_0Ir}{2\pi a^2}.
\]
磁场为绕 z 轴的方位向分量。

(3) 若 r>a，位移电流所包围的面积已经覆盖全部电场区域，所以 \(2\pi rB_\phi=\mu_0I\)，即 \(B_\phi(r)=\mu_0I/(2\pi r)\)。对同一闭合曲线选经过外导线的曲面则有传导电流 I，选穿过板间真空的曲面则有位移电流 I；两种曲面得到同一个磁场环量。若遗漏位移电流，安培定律会出现曲面依赖矛盾。

(4) 电容能量
\[
U=\frac{\epsilon_0}{2}E^2\pi a^2d,\qquad
\dot U=\epsilon_0E\dot E\pi a^2d=EId=VI.
\]
在侧面 r=a 处 \(B_\phi=\mu_0 I/(2\pi a)\)，坡印廷矢量 \(\boldsymbol S=\boldsymbol E\times\boldsymbol B/\mu_0\) 由 \(\hat z\times\hat\phi=-\hat r\) 指向电容内部。侧面总流入功率为 \(S\,2\pi ad=EI d=VI\)，与储能微分严格一致。这表明电磁能从板间侧方通过空间电磁场进入，而不是由真空中的真实传导电流带入。
